\section{System overview and management model}
\label{sec:system}

HMRS manages the state a search reads. The parsing service recovers document
content, the web application manages the collection and its files, and the
MCP server exposes these operations to agents \cite{mcp}. A document's
placement determines the scope in which it is visible. Its description
guides a catalogue-first search, and its index entries make passages
discoverable to lexical and dense search. To check these connections, HMRS
compares document files, the tree registry and per-library metadata after
selected management operations.

\subsection{Management levels and the tiered access path}
\label{sec:levels}

The management model has three levels: collections carry ownership and scope,
documents carry content and descriptive metadata, and segments carry the
text units used for passage retrieval or scoring. Four access views expose
those levels. The catalogue lists collections, the shelf lists documents
within a collection, the document view opens a body and the segment view
exposes passages and their source references. Table~\ref{tab:levels} links each
managed level to the operations and checks that exercise it.

\begin{table*}[t]
\centering
\small
\setlength{\tabcolsep}{4pt}
\renewcommand{\arraystretch}{1.15}
\caption{Management levels and the tiered access path. Each level names the
objects it maintains, the operations that write them, the state that describes
them, the tier that reads them, and the measurement that exercises it
(evidence labels E1--E9 denote the evaluation campaigns of
Section~\ref{sec:eval}; Table~S17 records their design). The selected access
policy determines which tiers are visited; a pinned
search can bypass catalogue navigation.}
\label{tab:levels}
\begin{tabular}{@{}>{\raggedright\arraybackslash}p{0.12\textwidth}>{\raggedright\arraybackslash}p{0.21\textwidth}>{\raggedright\arraybackslash}p{0.18\textwidth}>{\raggedright\arraybackslash}p{0.12\textwidth}>{\raggedright\arraybackslash}p{0.12\textwidth}>{\raggedright\arraybackslash}p{0.13\textwidth}@{}}
\toprule
Level & Objects and operations & State description & Read tier & Check & Evidence\\
\midrule
Collection & libraries and sub-libraries; agent assignment from parsed content; description and tag derivation; catalogue listing & library registry, per-library metadata, library-level index & catalogue & catalogue inventory and description round-trip & E1, E2, E3\\
Document & documents and their parts; duplicate review; move, rename, delete; forced re-index & document files, tree registry, per-library metadata & shelf & reconciliation on selected operations & E2\\
Segment & structural split; dense and lexical indices; optional evidence gate & chunk identifiers, per-segment scores with provenance & document and segment & index coverage and post-index probe & E1, E5--E9\\
\bottomrule
\end{tabular}
\end{table*}

The ingest agent uses catalogue descriptions to place parsed material and
derives document descriptions and tags from its content. Sub-libraries are
tree nodes with their own descriptions and indices. At the document level,
duplicate review identifies candidates and a lifecycle operation can move,
rename or delete a document. At the segment level, search returns indexed
chunks, while the optional evidence gate fetches candidate text and scores
query-conditioned passages with provenance.

Figure~\ref{fig:levels} uses the library analogy to show the transitions
between catalogue, shelf, document and segment views. A pinned search can
start inside selected collections and access indexed chunks directly.
Figure~\ref{fig:arch} shows how both paths reach the same document and
index services. The catalogue and shelf can be inspected independently:
an absent collection and an omitted document appear at different stages
of the access trace.

\begin{figure}[t]
\centering
\includegraphics[width=\columnwidth]{fig_library_analogy_vector}
\caption{Library analogy for a catalogue-first HMRS access path. The library catalogue, shelf, book and passage correspond to the collection catalogue, document listing, document body and segment. Evidence scoring is optional; pinned search may bypass the catalogue and shelf.}
\label{fig:levels}
\end{figure}

\setcounter{dbltopnumber}{1}
\begin{figure*}[t]
\centering
\includegraphics[width=\textwidth]{fig_service_architecture_vector}
\caption{HMRS service architecture. The browser and agent/CLI access the web and MCP interfaces. Parsing, management and retrieval operate on shared document files, the tree registry, library metadata and search indices. Laya scoring is optional, and the Neo4j graph supports navigation.}
\label{fig:arch}
\end{figure*}

Figure~\ref{fig:workflow} follows a document from acquisition and
organisation to evidence-based reporting. The ingest workflow governs
content review and placement. The query workflow determines whether to
consult the catalogue or search a pinned collection, and whether to use
Laya to score passages. The answer instructions distinguish a completed
review's searched scope and blind spots from an incomplete review caused
by tool failure.
Hybrid retrieval is also available as a tool and can be composed from
ordinary search calls. Experience retrieval and persona distillation reuse
the document infrastructure and define their own retrieval procedures.

Table~\ref{tab:functions} maps acquisition, organisation, lifecycle, query
and experience operations to the interface or workflow that checks them.
This allocation locates decisions at the appropriate boundary. For example,
the ingest agent chooses a collection, the save service validates the
destination identifier, and an audit checks the stored result.

\begin{table*}[t]
\centering
\small
\setlength{\tabcolsep}{4pt}
\caption{Management functions, the enforcement or workflow check associated
with each, and the evaluation that exercises it. The table separates API/MCP
controls from agent procedures and campaign audits; semantic choices such as
which collection owns a document are made and recorded by the workflow.
Evidence labels E1--E9 denote the campaigns in Section~\ref{sec:eval}.}
\label{tab:functions}
\begin{tabular}{@{}>{\raggedright\arraybackslash}p{0.13\textwidth}>{\raggedright\arraybackslash}p{0.37\textwidth}>{\raggedright\arraybackslash}p{0.30\textwidth}>{\raggedright\arraybackslash}p{0.11\textwidth}@{}}
\toprule
Group & Operations & Platform or workflow check & Evidence\\
\midrule
Acquisition & submit, parse (MinerU), quality review, duplicate review, structural split, store, index & parse-task state; workflow quality checks; entry-point-specific deduplication and indexing; post-index probe & E1--E2, S9\\
Organisation & create library or sub-library, content-informed placement, descriptions and tags, duplicate review & API validates identifiers and required fields; workflow checks descriptions and tags; audit compares files, tree and library metadata & E1--E2\\
Lifecycle & move, rename, delete, force re-index, graph build & operation-specific API updates; reconciliation audit detects selected stale or inconsistent state & E2\\
Query & tier selection, vector, keyword, catalogue and two-stage search, gate invocation & endpoint enforces tool allow/deny lists, timeouts and turn limits; invoked gate returns survivor provenance & E3, E5--E9\\
Experience & extract, approve, tier, search & experience tools apply their own retrieval, verification and credibility filters & E4\\
\bottomrule
\end{tabular}
\end{table*}

\begin{figure*}[t]
\centering
\includegraphics[width=\textwidth]{fig_workflow_vector}
\caption{Illustrative HMRS workflow from document acquisition to evidence-based reporting. Documents are parsed when needed, organised into collections and indexed according to the entry point. A question can use catalogue navigation or pinned search; optional Laya scoring precedes evidence inspection. A completed review without support is distinct from an incomplete search or tool failure, which may occur at any stage.}
\label{fig:workflow}
\end{figure*}

\subsection{Data model and object lifecycle}
\label{sec:datamodel}
\label{sec:lifecycle}

Documents are stored as Markdown with associated assets. A managed
collection's state spans the document files, \texttt{.tree-fs.json} and the
library's \texttt{.knowledge-base.yml}. The ingest workflow distinguishes
parsed, stored and indexed content and records completion after a retrieval
probe. The MCP save path can start background indexing; the web parsed-file
save path exposes indexing as a separate operation. Reconciliation compares
the three representations and identifies selected stale-index conditions.
A move, rename or re-index can therefore be checked through the resulting
file path, metadata and retrieval response.
